% GENERATED FILE. Edit canonical lessons, then run npm run generate:books.
% canonical-source-hash: fnv1a64:ade97be35d30cf18
% canonical-lessons: GU-C504-kumbhar, GU-C504-cakdo, GU-C504-garo, GU-C504-gudvu, GU-C504-kodiyu

\chapter{Clay and the Wheel}
\label{ch:gu-clay-and-the-wheel}
\hlchaptermodality{504}

\begin{chapteropening}
\textbf{By the end of this chapter:} \emph{I can talk about how a potter shapes clay.}

Complete the last lesson of chapter 504: I can talk about how a potter shapes clay.
\end{chapteropening}

\section[\texorpdfstring{\gu{કુંભાર} (kumbhār)}{કુંભાર (kumbhār)}]{\gu{કુંભાર} (kumbhār) — a potter}

\label{lesson:GU-C504-kumbhar}



\begin{quote}
Before the new one: say the Gujarati for a crack, then the Gujarati for to split (wood).
\end{quote}

\subsection*{You'll want to know: \gu{કુંભાર}}
\textbf{\gu{કુંભાર}} — \emph{kumbhār} — \textquotedblleft{}a potter\textquotedblright{}.

\begin{quote}
\gu{અમારા} \gu{ગામમાં} \gu{એક} \gu{કુંભાર} \gu{રહે} \gu{છે}.
\end{quote}

\begin{quote}
\emph{amārā gāmmā̃ ek kumbhār rahe che.}
\end{quote}

\begin{quote}
A potter lives in our village.
\end{quote}

A masculine noun.

\begin{cousinweb}
\textbf{\gu{કુંભાર}} (\emph{kumbhār}) is inherited from Sanskrit \emph{kumbhakāra}, \textquotedblleft{}a pot-maker\textquotedblright{}, built from \emph{kumbha}, \textquotedblleft{}a pot\textquotedblright{}, and \emph{kāra}, \textquotedblleft{}maker\textquotedblright{}; CDIAL lists the Gujarati word (CDIAL 3310).
\end{cousinweb}

\begin{grammarlens}[title={What you've built}]
One more word for describing what you lived and felt.
\end{grammarlens}

\subsection*{Your turn}
\begin{itemize}
\raggedright
  \item \emph{Say it:} \emph{kumbhār}
  \item \emph{Say it:} \emph{kumbhār}, once more
  \item \emph{Say it:} \emph{amārā gāmmā̃ ek kumbhār rahe che.} — A potter lives in our village.
  \item \emph{Recall:} say the Gujarati for a crack, then the Gujarati for to split (wood), then say \emph{kumbhār} again
\end{itemize}

\begin{tcolorbox}[breakable,colback=teal!4,colframe=teal!35!black,title={Before you move on}]
What does \gu{કુંભાર} mean? (\textquotedblleft{}A potter\textquotedblright{}.) Say it once more.
\end{tcolorbox}
\section[\texorpdfstring{\gu{ચાકડો} (cākḍo)}{ચાકડો (cākḍo)}]{\gu{ચાકડો} (cākḍo) — a potter's wheel}

\label{lesson:GU-C504-cakdo}



\begin{quote}
Before the new one: say the Gujarati for to split (wood), then the Gujarati for a potter.
\end{quote}

\subsection*{You'll want to know: \gu{ચાકડો}}
\textbf{\gu{ચાકડો}} — \emph{cākḍo} — \textquotedblleft{}a potter's wheel\textquotedblright{}.

\begin{quote}
\gu{કુંભાર} \gu{ચાકડા} \gu{પર} \gu{માટીનું} \gu{માટલું} \gu{બનાવે} \gu{છે}.
\end{quote}

\begin{quote}
\emph{kumbhār cākḍā par māṭīnũ māṭlũ banāve che.}
\end{quote}

\begin{quote}
The potter makes a clay pot on the wheel.
\end{quote}

A masculine noun.

\begin{cousinweb}
is built on \emph{cāk}, \textquotedblleft{}a wheel, a potter's wheel\textquotedblright{}, which CDIAL lists under Sanskrit \emph{cakrá}, \textquotedblleft{}a wheel\textquotedblright{} (CDIAL 4538); \textbf{\gu{ચાકડો}} adds the ending \emph{-ḍo}.
\end{cousinweb}

\begin{grammarlens}[title={What you've built}]
One more word for describing what you lived and felt.
\end{grammarlens}

\subsection*{Your turn}
\begin{itemize}
\raggedright
  \item \emph{Say it:} \emph{cākḍo}
  \item \emph{Say it:} \emph{cākḍo}, once more
  \item \emph{Say it:} \emph{kumbhār cākḍā par māṭīnũ māṭlũ banāve che.} — The potter makes a clay pot on the wheel.
  \item \emph{Recall:} say the Gujarati for to split (wood), then the Gujarati for a potter, then say \emph{cākḍo} again
\end{itemize}

\begin{tcolorbox}[breakable,colback=teal!4,colframe=teal!35!black,title={Before you move on}]
What does \gu{ચાકડો} mean? (\textquotedblleft{}A potter's wheel\textquotedblright{}.) Say it once more.
\end{tcolorbox}
\section[\texorpdfstring{\gu{ગારો} (gāro)}{ગારો (gāro)}]{\gu{ગારો} (gāro) — mud plaster}

\label{lesson:GU-C504-garo}



\begin{quote}
Before the new one: say the Gujarati for a potter, then the Gujarati for a potter's wheel.
\end{quote}

\subsection*{You'll want to know: \gu{ગારો}}
\textbf{\gu{ગારો}} — \emph{gāro} — \textquotedblleft{}mud plaster\textquotedblright{}.

\begin{quote}
\gu{અમારા} \gu{ગામનાં} \gu{જૂનાં} \gu{ઘર} \gu{ગારાનાં} \gu{છે}.
\end{quote}

\begin{quote}
\emph{amārā gāmnā̃ jūnā̃ ghar gārānā̃ che.}
\end{quote}

\begin{quote}
The old houses in our village are made of mud.
\end{quote}

A masculine noun.

\begin{cousinweb}
\textbf{\gu{ગારો}} (\emph{gāro}) comes from \emph{*gāra}, a form CDIAL reconstructs, and CDIAL lists the Gujarati word as earth and water mixed for building a wall (CDIAL 4137).
\end{cousinweb}

\begin{grammarlens}[title={What you've built}]
One more word for describing what you lived and felt.
\end{grammarlens}

\subsection*{Your turn}
\begin{itemize}
\raggedright
  \item \emph{Say it:} \emph{gāro}
  \item \emph{Say it:} \emph{gāro}, once more
  \item \emph{Say it:} \emph{amārā gāmnā̃ jūnā̃ ghar gārānā̃ che.} — The old houses in our village are made of mud.
  \item \emph{Recall:} say the Gujarati for a potter, then the Gujarati for a potter's wheel, then say \emph{gāro} again
\end{itemize}

\begin{tcolorbox}[breakable,colback=teal!4,colframe=teal!35!black,title={Before you move on}]
What does \gu{ગારો} mean? (\textquotedblleft{}Mud plaster\textquotedblright{}.) Say it once more.
\end{tcolorbox}
\section[\texorpdfstring{\gu{ગૂંદવું} (gū̃dvũ)}{ગૂંદવું (gū̃dvũ)}]{\gu{ગૂંદવું} (gū̃dvũ) — to knead clay}

\label{lesson:GU-C504-gudvu}



\begin{quote}
Before the new one: say the Gujarati for a potter's wheel, then the Gujarati for mud plaster.
\end{quote}

\subsection*{You'll want to know: \gu{ગૂંદવું}}
\textbf{\gu{ગૂંદવું}} — \emph{gū̃dvũ} — \textquotedblleft{}to knead clay\textquotedblright{}.

\begin{quote}
\gu{કુંભારે} \gu{પગથી} \gu{માટી} \gu{ગૂંદી}.
\end{quote}

\begin{quote}
\emph{kumbhāre pagthī māṭī gū̃dī.}
\end{quote}

\begin{quote}
The potter kneaded the clay with his feet.
\end{quote}

Past, with \gu{તેણે} (he or she …): \gu{તેણે} \gu{માટી} \gu{ગૂંદી}, \emph{teṇe māṭī gū̃dī}.

\begin{cousinweb}
\textbf{\gu{ગૂંદવું}} (\emph{gū̃dvũ}) is said of clay or of a stiff dough worked long and hard; the past agrees with what is kneaded: \textbf{\gu{માટી} \gu{ગૂંદી}}. CDIAL lists Hindi \emph{gū̃dhnā}, \textquotedblleft{}to plait, mix, knead\textquotedblright{}, under Sanskrit \emph{gupháti}, \textquotedblleft{}strings together\textquotedblright{} (CDIAL 4205), beside a Gujarati verb for plaiting; it gives no Gujarati form for kneading, so its history is not settled.
\end{cousinweb}

\begin{grammarlens}[title={What you've built}]
One more word for describing what you lived and felt.
\end{grammarlens}

\subsection*{Your turn}
\begin{itemize}
\raggedright
  \item \emph{Say it:} \emph{gū̃dvũ}
  \item \emph{Say it:} \emph{gū̃dvũ}, once more
  \item \emph{Say it:} \emph{kumbhāre pagthī māṭī gū̃dī.} — The potter kneaded the clay with his feet.
  \item \emph{Recall:} say the Gujarati for a potter's wheel, then the Gujarati for mud plaster, then say \emph{gū̃dvũ} again
\end{itemize}

\begin{tcolorbox}[breakable,colback=teal!4,colframe=teal!35!black,title={Before you move on}]
What does \gu{ગૂંદવું} mean? (\textquotedblleft{}To knead clay\textquotedblright{}.) Say it once more.
\end{tcolorbox}
\section[\texorpdfstring{\gu{કોડિયું} (koḍiyũ)}{કોડિયું (koḍiyũ)}]{\gu{કોડિયું} (koḍiyũ) — a clay lamp}

\label{lesson:GU-C504-kodiyu}



\begin{quote}
Before the new one: say the Gujarati for mud plaster, then the Gujarati for to knead clay.
\end{quote}

\subsection*{You'll want to know: \gu{કોડિયું}}
\textbf{\gu{કોડિયું}} — \emph{koḍiyũ} — \textquotedblleft{}a clay lamp\textquotedblright{}.

\begin{quote}
\gu{દિવાળીમાં} \gu{અમે} \gu{બારીમાં} \gu{કોડિયાં} \gu{મૂકીએ} \gu{છીએ}.
\end{quote}

\begin{quote}
\emph{divāḷīmā̃ ame bārīmā̃ koḍiyā̃ mūkīe chīe.}
\end{quote}

\begin{quote}
At Diwali we put clay lamps in the windows.
\end{quote}

A neuter noun.

\begin{cousinweb}
\textbf{\gu{કોડિયું}} (\emph{koḍiyũ}) is filed by CDIAL under Sanskrit \emph{kuṭa}, \textquotedblleft{}a water-pot, a pitcher\textquotedblright{}, and CDIAL lists the Gujarati word as an earthen cup for oil and a wick (CDIAL 3227). The plural is \textbf{\gu{કોડિયાં}}.
\end{cousinweb}

\begin{grammarlens}[title={What you've built}]
One more word for describing what you lived and felt.
\end{grammarlens}

\subsection*{Your turn}
\begin{itemize}
\raggedright
  \item \emph{Say it:} \emph{koḍiyũ}
  \item \emph{Say it:} \emph{koḍiyũ}, once more
  \item \emph{Say it:} \emph{divāḷīmā̃ ame bārīmā̃ koḍiyā̃ mūkīe chīe.} — At Diwali we put clay lamps in the windows.
  \item \emph{Recall:} say the Gujarati for mud plaster, then the Gujarati for to knead clay, then say \emph{koḍiyũ} again
\end{itemize}

\begin{tcolorbox}[breakable,colback=teal!4,colframe=teal!35!black,title={Before you move on}]
What does \gu{કોડિયું} mean? (\textquotedblleft{}A clay lamp\textquotedblright{}.) Say it once more.
\end{tcolorbox}
